%======================================================================
\chapter{Mullikin's theorem and the Jordan curve theorem}\label{ch:jordan}
%======================================================================

This chapter proves Moore's Theorem~11: \emph{a simple closed curve of elements has exactly two
complementary domains, one bounded and one unbounded, and is the boundary of each}. Applied to the
decomposition of the plane into points, it is the Jordan curve theorem. The plane input is a theorem
of A.~M.~Mullikin (1922) on the union of two non-separating continua, which we prove with
logarithms, as in Volume~4, instead of Mullikin's polygons.

Points of the plane are identified with complex numbers (\vfour{Section~0.4}). Recall from Volume~4:
two logarithms of the same function on a connected set differ by a constant in $2\pi i\Z$
(\vfour{Lemma~1.3}); if $y$ is in the component of $x$ in the complement of a closed set $K$, then
$(z-x)/(z-y)$ has a logarithm on $K$ (\vfour{Theorem~2.4}); if that component is bounded, then for
every integer $a\ne0$ and every $u$ continuous and nowhere zero on its closure, $(z-x)^au(z)$ has no
logarithm on $K$ (\vfour{Theorem~2.5}); and if $K$ is compact with connected complement, every
continuous nowhere zero function on $K$ has a logarithm on $K$ (Borsuk, \vfour{Theorem~2.10}).

\section{Mullikin's theorem}

\begin{theorem}[Mullikin 1922, Theorem~5, corrected]\label{thm:mullikin}
Let $M_1,M_2\subseteq\R^2$ be continua neither of which separates the plane, and suppose
$M_1\cap M_2=K_1\cup K_2$ with $K_1,K_2$ disjoint, closed and connected. Then $\R^2\setminus(M_1\cup
M_2)$ has exactly two components: it is the union of two disjoint connected open sets.
\end{theorem}

The theorem has two halves, proved separately.

\begin{proposition}[At least two components]\label{prop:mullikin-two}
Let $M_1,M_2$ be continua with $M_1\cap M_2=K_1\cup K_2$, where $K_1,K_2$ are disjoint, nonempty and
closed. Then $M_1\cup M_2$ separates the plane.
\end{proposition}

\begin{proof}
Suppose $\R^2\setminus(M_1\cup M_2)$ is connected. By Urysohn's lemma there is a continuous
$\varphi\colon\R^2\to[0,1]$ with $\varphi=0$ on $K_1$ and $\varphi=1$ on $K_2$. Define $f$ on $M_1\cup
M_2$ by $f=1$ on $M_1$ and $f=e^{2\pi i\varphi}$ on $M_2$; on $M_1\cap M_2$ both formulas give $1$, and
$M_1,M_2$ are closed, so $f$ is continuous and nowhere zero. By Borsuk's theorem $f$ has a logarithm
$g$ on the compact set $M_1\cup M_2$. Take $k_1\in K_1$, $k_2\in K_2$. On the connected set $M_1$, $g$ and
the constant $0$ are logarithms of $1$, so $g(k_1)=g(k_2)$. On the connected set $M_2$, $g$ and
$2\pi i\varphi$ are logarithms of $e^{2\pi i\varphi}$, so $g(k_2)-g(k_1)=2\pi i(\varphi(k_2)-\varphi(k_1))=2\pi
i$. Contradiction.
\end{proof}

\formal{\lean{HPC.Moore.not_isConnected_compl_union}.}

\begin{proposition}[At most two components; Eilenberg]\label{prop:mullikin-three}
Let $M_1,M_2$ be compact sets with connected complements, and $M_1\cap M_2=K_1\cup K_2$ with $K_1,K_2$
connected. Then $\R^2\setminus(M_1\cup M_2)$ has at most two components.
\end{proposition}

\begin{proof}
Put $M=M_1\cup M_2$, a compact set. Suppose $M$ had three components $C_p,C_q,C_y$ of its complement,
containing $p,q,y$. At most one component is unbounded (\vfour{Lemma~2.3}); renaming, we may assume
$C_p$ and $C_q$ are bounded.

For $x\in\set{p,q}$ and $i\in\set{1,2}$, the point $x$ and $y$ lie off $M_i$, whose complement is
connected; so $h_x(z)=(z-x)/(z-y)$ has a logarithm $g_{ix}$ on $M_i$ (\vfour{Theorem~2.4}). On each
connected set $K_j$ the two logarithms $g_{1x},g_{2x}$ of $h_x$ differ by a constant: $g_{1x}-g_{2x}=2\pi
i\,n_{jx}$ on $K_j$ with $n_{jx}\in\Z$ (if $K_j$ is empty there is nothing to say; take $n_{jx}=0$).

Choose integers $a,b$, not both zero, with $a(n_{1p}-n_{2p})+b(n_{1q}-n_{2q})=0$ (if both differences
vanish take $a=1,b=0$; otherwise $a=n_{1q}-n_{2q}$, $b=-(n_{1p}-n_{2p})$). Define $G$ on $M$ by
\[
G=a\,g_{1p}+b\,g_{1q}\ \text{on }M_1,\qquad G=a\,g_{2p}+b\,g_{2q}+2\pi i(a\,n_{1p}+b\,n_{1q})\ \text{on }M_2\setminus M_1.
\]
On $K_1$ the two expressions agree by definition of the $n_{1x}$, and on $K_2$ they agree by the choice
of $a,b$; so the second expression, extended to all of $M_2$, agrees with the first on $M_1\cap M_2$,
and $G$ is continuous on $M$ ($M_1,M_2$ closed). Moreover $e^{G}=h_p^a\,h_q^b$ on $M$.

Say $a\neq0$ (the case $b\ne0$ is symmetric, exchanging $p$ and $q$). Then
$h_p^ah_q^b=(z-p)^a\,u(z)$ with $u(z)=(z-y)^{-a}h_q(z)^b$. The closure of $C_p$ lies in $C_p\cup M$,
which contains neither $y$ nor $q$, so $u$ is continuous and nowhere zero on $\cl{C_p}$. By
\vfour{Theorem~2.5}, $(z-p)^au(z)$ has no logarithm on $M$; but $G$ is one. Contradiction.
\end{proof}

\formal{\lean{HPC.Moore.exists_jump}, \lean{HPC.Moore.no_log_pow_of_bounded_component},
\lean{HPC.Moore.not_two_bounded_components}, \lean{HPC.Moore.compl_union_two_components}.}

\begin{proof}[Proof of Theorem~\ref{thm:mullikin}]
$K_1,K_2$ are nonempty (connected sets are nonempty), so the complement of $M_1\cup M_2$ has at least
two components (Proposition~\ref{prop:mullikin-two}) and at most two
(Proposition~\ref{prop:mullikin-three}). The components are open, since $M_1\cup M_2$ is closed.
\end{proof}

\formal{\lean{HPC.Moore.mullikin_theorem5}.}

\begin{remark}\label{rem:mullikin-printed}
\leavevmode
\begin{enumerate}
\item Mullikin states the theorem with $K_1,K_2$ merely mutually exclusive connected sets. This is
false: for $M_1=M_2$ a segment, $M_1\cap M_2$ is the union of two disjoint connected (not closed)
pieces, and the complement is connected (\lean{HPC.Moore.not_mullikin_theorem5_printed}). Moore uses
the theorem with $K_1,K_2$ two elements, which are closed.
\item Mullikin's Theorem~4 says that the union of two plane continua, neither separating the plane,
separates the plane if and only if their intersection is not connected. With her convention that the
empty set is connected, the case $M_1\cap M_2=\emptyset$ says that the union of two disjoint
non-separating continua does not separate the plane; this is true (Janiszewski's theorem with empty
intersection, \lean{HPC.Moore.mullikin_theorem4_disjoint}). With the modern convention (the empty set
is not connected) the same sentence would be false for two points
(\lean{HPC.Moore.not_mullikin_theorem4_disjoint_mathlib}).
\end{enumerate}
\end{remark}

\section{Theorem 11}

We return to the standing hypothesis of Chapter~\ref{ch:usc}.

\begin{lemma}\label{lem:frontier-scc}
Let $J$ be a simple closed curve of elements and $I\subseteq G$ a nonempty open set with $I\cap
J=\emptyset$ and $\fr I\subseteq J$. Then $\fr I=J$.
\end{lemma}

\begin{proof}
Put $B=\fr I$, a closed subset of $J$, and suppose $x\in J\setminus B$. Since $J$ is homeomorphic to
a circle (Proposition~\ref{prop:arc-elt-top}(c)) and $B$ is closed, there are elements $p_1\ne q_1$ of
$J$ such that one of the two arcs of $J$ from $p_1$ to $q_1$ (Theorem~\ref{thm:m25-5-10}(e)), call it
$C$, contains $x$ as an interior point and misses $B$. Let $A$ be the other arc; then $B\subseteq A$ and
$x\notin A$. By Theorem~\ref{thm:m25-9}, $G\setminus A$ is connected. Now $I\subseteq G\setminus A$, $I$
is open, and its closure $I\cup B$ meets $G\setminus A$ only in $I$; so $I$ is a nonempty open and
closed subset of $G\setminus A$, hence $I=G\setminus A$. But $x\in G\setminus A$ and $x\in J$, while
$I\cap J=\emptyset$.
\end{proof}

\formal{\lean{HPC.Moore.frontier_eq_of_scc}.}

\begin{theorem}[Moore, Theorem~11]\label{thm:m25-11}
Let $J$ be a simple closed curve of elements. Then $G\setminus J=I\cup E$ with $I,E$ disjoint,
open and connected, $I$ bounded, $E$ unbounded, and $\fr I=\fr E=J$.
\end{theorem}

\begin{proof}
Choose distinct $p,q\in J$ and, by Theorem~\ref{thm:m25-5-10}(e), arcs $a,b$ of elements from $p$
to $q$ with $a\cup b=J$ and $a\cap b=\set{p,q}$. By Theorem~\ref{thm:m25-9} and
Corollary~\ref{cor:pts-nonsep}, neither $\pts a$ nor $\pts b$ separates the plane. They are continua
(compact by Proposition~\ref{prop:arc-elt-top}, connected by Proposition~\ref{prop:pts-conn}(a)), and
$\pts a\cap\pts b=\pts{\set{p,q}}=p\cup q$ is the union of the two disjoint continua $p$ and $q$
(elements are compact and connected). By Mullikin's theorem, $\R^2\setminus\pts J=U\cup V$ with
$U,V$ disjoint, connected and open.

Each element $g\notin J$ is connected and lies in $U\cup V$, so it lies in $U$ or in $V$; hence $U$
and $V$ are unions of elements. Put $I'=\pi(U)$, $E'=\pi(V)$: they are open (their point sets are $U,V$),
connected, disjoint, and $I'\cup E'=G\setminus J$. As $J$ is bounded, exactly one component of
$G\setminus J$ is unbounded (Corollary~\ref{cor:quot-tools}(c)), and not both are bounded (else $G=I'\cup
E'\cup J$ would be bounded). Let $I$ be the bounded one and $E$ the unbounded one.

$I$ is open, and it is closed in $G\setminus J$ (its complement there is the open set $E$); so $\cl
I\subseteq I\cup J$ and $\fr I\subseteq J$. Lemma~\ref{lem:frontier-scc} gives $\fr I=J$; likewise $\fr
E=J$.
\end{proof}

\formal{\lean{HPC.Moore.m25_theorem11}, \lean{HPC.Moore.sccInterior_spec}.}

\begin{remark}
Moore's proof says that $J$ is the boundary of each domain because ``it easily follows, with the help
of Theorems~10 and~6, that $B$ is a simple continuous arc''. Theorem~6 needs $B=\fr I$ to be
connected, which is not known at that point. Lemma~\ref{lem:frontier-scc} replaces this step; it uses
Theorems~9 and~10 only.
\end{remark}

The bounded domain $I$ is the \emph{interior} of $J$ and $E$ its \emph{exterior}.

\section{The Jordan curve theorem}

\begin{theorem}[Jordan curve theorem]\label{thm:jct}
Let $J\subseteq\R^2$ be homeomorphic to a circle. Then $\R^2\setminus J=I\cup E$ with $I,E$ disjoint,
open and connected, $I$ bounded and $E$ unbounded, and $\fr I=\fr E=J$.
\end{theorem}

\begin{proof}
Apply Theorem~\ref{thm:m25-11} to the decomposition of the plane into points
(Section~\ref{sec:points}): there $G$ is the plane itself, and by
Proposition~\ref{prop:arc-elt-top}(c) the simple closed curves of elements are exactly the subsets
of the plane homeomorphic to a circle.
\end{proof}

\formal{\lean{HPC.Jones.jordan_curve_theorem}.}

The proof uses: logarithms and Borsuk's theorem (Volume~4), Brouwer's theorem and Janiszewski's
theorem (Volume~4), the characterizations of arcs and circles (Chapter~\ref{ch:arcs}), Theorems~2,
6, 9 and~10 (Chapters~\ref{ch:usc} and~\ref{ch:elements}) and Mullikin's theorem. It does not use
Moore's characterization of the plane (Chapters~\ref{ch:axioms}--\ref{ch:plane}). In the decomposition
into points, Theorem~2 and Theorem~9 are statements about plane continua and plane arcs, and the
whole argument is a proof of the Jordan curve theorem by Janiszewski's and Borsuk's theorems, in the
spirit of Eilenberg's.

\section*{Exercises}
\addcontentsline{toc}{section}{Exercises}

\begin{exercise}
Deduce from Proposition~\ref{prop:mullikin-two} that a simple closed curve in the plane separates the
plane, using only that an arc does not separate the plane.
\end{exercise}

\begin{exercise}
Generalize Proposition~\ref{prop:mullikin-three}: if $M_1\cap M_2$ has $n$ components, the complement of
$M_1\cup M_2$ has at most $n$ components.
\end{exercise}

\begin{exercise}
Show that the theta-curve (a circle together with a diameter) has exactly three complementary
domains, using Theorem~\ref{thm:jct} and Proposition~\ref{prop:mullikin-three}.
\end{exercise}
